% Combined source bibliography, rebuilt 2026-09-29 from the current paper files.
% sys:N = cwaight_systems_paper.tex ref N; sep:N = Draft_11.tex ref N;
% idc:N (bib_idetc.tex) = IDETC 2025 ref N. build_bib.py dedupes by title
% and renumbers to integer keys in order of first citation.

\bibitem{sys:1} C. Kitts and M. Egerstedt, "Design, control, and applications of real-world multirobot systems [from the guest editors]," \textit{IEEE Robotics \& Automation Magazine}, vol. 15, no. 1, p. 8, Mar. 2008.
\bibitem{sys:2} J. Cortés and M. Egerstedt, "Coordinated control of multi-robot systems: A survey," \textit{SICE Journal of Control, Measurement, and System Integration}, vol. 10, no. 6, pp. 495--503, 2017.
\bibitem{sys:3} P. Ögren, E. Fiorelli, and N. E. Leonard, "Cooperative control of mobile sensor networks: Adaptive gradient climbing in a distributed environment," \textit{IEEE Transactions on Automatic Control}, vol. 49, no. 8, pp. 1292--1302, 2004.
\bibitem{sys:4} L. Briñón-Arranz, A. Renzaglia, and L. Schenato, "Multirobot symmetric formations for gradient and hessian estimation with application to source seeking," \textit{IEEE Transactions on Robotics}, vol. 35, no. 3, pp. 782--789, 2019.
\bibitem{sys:5} C. A. Kitts, R. T. McDonald, and M. A. Neumann, "Adaptive navigation control primitives for multirobot clusters: Extrema finding, contour following, ridge/trench following, and saddle point station keeping," \textit{IEEE Access}, vol. 6, pp. 17625--17642, 2018.
\bibitem{sys:6} R. T. McDonald, C. A. Kitts, and M. A. Neumann, "Experimental implementation and verification of scalar field ridge, trench, and saddle point maneuvers using multirobot adaptive navigation," \textit{IEEE Access}, vol. 7, pp. 62950--62961, 2019.
\bibitem{sys:7} R. K. Lee et al., "Multiple UAV adaptive navigation for three-dimensional scalar fields," \textit{IEEE Access}, vol. 9, pp. 122626--122654, 2021.
\bibitem{sys:8} R. T. McDonald, C. A. Kitts, and M. A. Neumann, "Navigation of scalar fronts with multirobot clusters in simulation and experiment," \textit{IEEE Systems Journal}, vol. 14, no. 3, pp. 3755--3766, 2020.
\bibitem{sys:9} R. K. Lee, C. A. Kitts, M. A. Neumann, and R. T. McDonald, "3-d adaptive navigation: Multirobot formation control for seeking and tracking of a moving source," in \textit{Proc. IEEE Aerospace Conference}, pp. 1--12, 2020.
\bibitem{sys:10} T. Adamek, C. A. Kitts, and I. Mas, "Gradient-based cluster space navigation for autonomous surface vessels," \textit{IEEE/ASME Transactions on Mechatronics}, vol. 20, no. 2, pp. 506--518, 2014.
\bibitem{sys:11} F. Pagano, N. De Carli, E. Restrepo, A. Marino, and P. Robuffo Giordano, "Distributed multi-robot active-sensing of a diffusive source," \textit{IEEE Robotics and Automation Letters}, vol. 10, no. 10, pp. 10807--10814, 2025.
\bibitem{sys:13} S. Li, R. Kong, and Y. Guo, "Cooperative distributed source seeking by multiple robots: Algorithms and experiments," \textit{IEEE/ASME Transactions on Mechatronics}, vol. 19, no. 6, pp. 1810--1820, 2014.
\bibitem{sys:14} Z. Deng and J. Luo, "Fully distributed algorithms for constrained nonsmooth optimization problems of general linear multiagent systems and their application," \textit{IEEE Transactions on Automatic Control}, vol. 69, no. 2, pp. 1377--1384, 2023.
\bibitem{sys:15} W. Li, J. A. Farrell, S. Pang, and R. M. Arrieta, "Moth-inspired chemical plume tracing on an autonomous underwater vehicle," \textit{IEEE Transactions on Robotics}, vol. 22, no. 2, pp. 292--307, 2006.
\bibitem{sys:18} S. R. Ramp et al., "Preparing to predict: the second autonomous ocean sampling network (AOSN-II) experiment in the Monterey Bay," \textit{Deep Sea Research Part II: Topical Studies in Oceanography}, vol. 56, no. 3-5, pp. 68--86, 2009.
\bibitem{sys:19} J. Das et al., "Coordinated sampling of dynamic oceanographic features with underwater vehicles and drifters," \textit{The International Journal of Robotics Research}, vol. 31, no. 5, pp. 626--646, 2012.
\bibitem{sys:20} J. Wang et al., "Dynamic plume tracking by cooperative robots," \textit{IEEE/ASME Transactions on Mechatronics}, vol. 24, no. 2, pp. 609--620, 2019.
\bibitem{sys:21} K. Garg and D. Panagou, "Finite-time estimation and control for multi-aircraft systems under wind and dynamic obstacles," \textit{Journal of Guidance, Control, and Dynamics}, vol. 42, no. 7, pp. 1489--1505, 2019.
\bibitem{sys:22} C. Waight and C. Kitts, "A functional indoor testbed for multirobot adaptive navigation in vector field environments," in \textit{Proc. ASME Int. Design Engineering Technical Conf. Computers and Information in Engineering Conf.}, vol. 89251, 2025.
\bibitem{sys:23} M. Michini, M. A. Hsieh, E. Forgoston, and I. B. Schwartz, "Robotic tracking of coherent structures in flows," \textit{IEEE Transactions on Robotics}, vol. 30, no. 3, pp. 593--603, 2014.
\bibitem{sys:24} D. R. Nelson, D. B. Barber, T. W. McLain, and R. W. Beard, "Vector field path following for miniature air vehicles," \textit{IEEE Transactions on Robotics}, vol. 23, no. 3, pp. 519--529, 2007.
\bibitem{sys:25} W. Yao, H. G. de Marina, B. Lin, and M. Cao, "Singularity-free guiding vector field for robot navigation," \textit{IEEE Transactions on Robotics}, vol. 37, no. 4, pp. 1206--1221, 2021.
\bibitem{sys:26} D. Panagou, "Motion planning and collision avoidance using navigation vector fields," in \textit{Proc. IEEE International Conference on Robotics and Automation}, 2014, pp. 2513--2518.
\bibitem{sys:28} X. Wang et al., "Adaptive vector field guidance without a priori knowledge of course dynamics and wind," \textit{IEEE/ASME Transactions on Mechatronics}, vol. 27, no. 6, pp. 4597--4607, 2022.
\bibitem{sys:29} Y. Zhang, M. A. Godin, J. G. Bellingham, and J. P. Ryan, "Using an autonomous underwater vehicle to track a coastal upwelling front," \textit{IEEE Journal of Oceanic Engineering}, vol. 37, no. 3, pp. 338--347, 2012.
\bibitem{sys:32} D. Kularatne and M. A. Hsieh, "Tracking attracting manifolds in flows," \textit{Autonomous Robots}, vol. 41, no. 8, pp. 1575--1588, 2017.
\bibitem{sys:36} X. Tricoche, G. Scheuermann, and H. Hagen, "Continuous topology simplification of planar vector fields," in \textit{Proc. IEEE Visualization}, pp. 159--166, 2001.
\bibitem{sys:37} W. G. Pichel et al., "Marine debris collects within the North Pacific subtropical convergence zone," \textit{Marine Pollution Bulletin}, vol. 54, no. 8, pp. 1207--1211, 2007.
\bibitem{sys:38} M. Mateus, R. Canelas, L. Pinto, and N. Vaz, "When tragedy strikes: potential contributions from ocean observation to search and rescue operations after drowning accidents," \textit{Frontiers in Marine Science}, vol. 7, p. 55, 2020.
\bibitem{sys:40} P. Keramea, K. Spanoudaki, G. Zodiatis, G. Gikas, and G. Sylaios, "Oil spill modeling: A critical review on current trends, perspectives, and challenges," \textit{Journal of Marine Science and Engineering}, vol. 9, no. 2, p. 181, 2021.
\bibitem{sys:41} J. Wang, Y. Lin, R. Liu, and J. Fu, "Odor source localization of multi-robots with swarm intelligence algorithms: A review," \textit{Frontiers in Neurorobotics}, vol. 16, p. 949888, 2022.
\bibitem{sys:43} C. Dong et al., "Oceanic mesoscale eddies," \textit{Ocean-Land-Atmosphere Research}, vol. 4, p. 0081, 2025.
\bibitem{sys:44} Z. Li, X. Wang, J. Hu, F. P. Andutta, and Z. Liu, "A study on an anticyclonic-cyclonic eddy pair off Fraser Island, Australia," \textit{Frontiers in Marine Science}, vol. 7, p. 594358, 2020.
\bibitem{sys:45} R. Deogharia, H. Gupta, S. Sil, A. Gangopadhyay, and A. Shee, "On the evidence of helico-spiralling recirculation within coherent cores of eddies using Lagrangian approach," \textit{Scientific Reports}, vol. 14, no. 1, p. 11014, 2024.
\bibitem{sys:30} T. Inanc, S. C. Shadden, and J. E. Marsden, "Optimal trajectory generation in ocean flows," in \textit{Proc. American Control Conference}, pp. 674--679, 2005.
\bibitem{sys:31} R. N. Smith, M. Schwager, S. L. Smith, B. H. Jones, D. Rus, and G. S. Sukhatme, "Persistent ocean monitoring with underwater gliders: Adapting sampling resolution," \textit{Journal of Field Robotics}, vol. 28, no. 5, pp. 714--741, 2011.
\bibitem{sys:33} G. Knizhnik, P. Li, X. Yu, and M. A. Hsieh, "Flow-based control of marine robots in gyre-like environments," in \textit{Proc. IEEE International Conference on Robotics and Automation (ICRA)}, pp. 3047--3053, 2022.
\bibitem{sys:46} Y. Jiao, H. Hang, J. Merel, and E. Kanso, "Sensing flow gradients is necessary for learning autonomous underwater navigation," \textit{Nature Communications}, vol. 16, no. 1, p. 3044, 2025.
\bibitem{sys:47} D. Rajabhandharaks, R. T. McDonald, M. A. Neumann, and C. A. Kitts, "Indoor testbed for vector field multirobot adaptive navigation: Feasibility study," in \textit{Proc. ASME Int. Design Engineering Technical Conf. Computers and Information in Engineering Conf.}, vol. 59292, p. V009T12A043, 2019.
\bibitem{sys:48} R. Cooper et al., "Initial Study of Multirobot Adaptive Navigation for Exploring Environmental Vector Fields," in \textit{Proc. IEEE Aerospace Conference}, 2019.
\bibitem{sys:49} Z. Zhou, J. Liu, and J. Yu, "A survey of underwater multi-robot systems," \textit{IEEE/CAA Journal of Automatica Sinica}, vol. 9, no. 1, pp. 1--18, 2021.
\bibitem{sys:50} C. A. Kitts and I. Mas, "Cluster space specification and control of mobile multirobot systems," \textit{IEEE/ASME Transactions on Mechatronics}, vol. 14, no. 2, pp. 207--218, 2009.
\bibitem{sys:51} I. Mas and C. Kitts, "Obstacle avoidance policies for cluster space control of nonholonomic multirobot systems," \textit{IEEE/ASME Transactions on Mechatronics}, vol. 17, no. 6, pp. 1068--1079, 2012.
\bibitem{sys:52} S. Tomer et al., "A low-cost indoor testbed for multirobot adaptive navigation research," in \textit{Proc. IEEE Aerospace Conference}, pp. 1--12, 2018.
\bibitem{sys:53} S. T. Hart and C. A. Kitts, "Unifying control architecture for reactive particle swarms," \textit{IEEE/ASME Transactions on Mechatronics}, vol. 28, no. 2, pp. 873--883, 2022.

\bibitem{sep:1} J. Cortés and M. Egerstedt, ``Coordinated control of multi-robot systems: A survey,'' \textit{SICE J. Control Meas. Syst. Integr.}, vol. 10, no. 6, pp. 495--503, 2017.
\bibitem{sep:2} T. Adamek, C. A. Kitts, and I. Mas, ``Gradient-based cluster space navigation for autonomous surface vessels,'' \textit{IEEE/ASME Trans. Mechatronics}, vol. 20, no. 2, pp. 506--518, 2015.
\bibitem{sep:3} C. A. Kitts, R. T. McDonald, and M. A. Neumann, ``Adaptive navigation control primitives for multirobot clusters: Extrema finding, contour following, ridge/trench following, and saddle point station keeping,'' \textit{IEEE Access}, vol. 6, pp. 17625--17642, 2018.
\bibitem{sep:4} R. T. McDonald et al., ``Navigation of scalar fronts with multirobot clusters in simulation and experiment,'' \textit{IEEE Syst. J.}, vol. 14, no. 3, pp. 3755--3766, 2020.
\bibitem{sep:5} L. Briñón-Arranz, A. Renzaglia, and L. Schenato, ``Multirobot symmetric formations for gradient and Hessian estimation with application to source seeking,'' \textit{IEEE Trans. Robot.}, vol. 35, no. 3, pp. 782--789, 2019.
\bibitem{sep:6} P. Ögren, E. Fiorelli, and N. E. Leonard, ``Cooperative control of mobile sensor networks: Adaptive gradient climbing in a distributed environment,'' \textit{IEEE Trans. Autom. Control}, vol. 49, no. 8, pp. 1292--1302, 2004.
\bibitem{sep:7} R. T. McDonald, C. A. Kitts, and M. A. Neumann, ``Experimental implementation and verification of scalar field ridge, trench, and saddle point maneuvers using multirobot adaptive navigation,'' \textit{IEEE Access}, vol. 7, pp. 62950--62961, 2019.
\bibitem{sep:8} A. Okubo, ``Horizontal dispersion of floatable particles in the vicinity of velocity singularities such as convergences,'' \textit{Deep-Sea Res.}, vol. 17, no. 3, pp. 445--454, 1970.
\bibitem{sep:9} P. Keramea, K. Spanoudaki, G. Zodiatis, G. Gikas, and G. Sylaios, ``Oil spill modeling: A critical review on current trends, downscaling, and applications,'' \textit{J. Mar. Sci. Eng.}, vol. 9, no. 2, art. 181, 2021.
\bibitem{sep:10} C. Waight and C. A. Kitts, ``Adaptive navigation of multirobot systems to critical points in 2D vector fields,'' submitted for publication.
\bibitem{sep:11} T. Peacock and G. Haller, ``Lagrangian coherent structures: The hidden skeleton of fluid flows,'' \textit{Phys. Today}, vol. 66, no. 2, pp. 41--47, 2013.
\bibitem{sep:12} M. Serra et al., ``Search and rescue at sea aided by hidden flow structures,'' \textit{Nat. Commun.}, vol. 11, no. 1, p. 2525, 2020.
\bibitem{sep:13} J. Das et al., ``Coordinated sampling of dynamic oceanographic features with underwater vehicles and drifters,'' \textit{Int. J. Robot. Res.}, vol. 31, no. 5, pp. 626--646, 2012.
\bibitem{sep:14} S. McCammon et al., ``Ocean front detection and tracking using a team of heterogeneous marine vehicles,'' \textit{J. Field Robot.}, vol. 38, no. 6, pp. 854--881, 2021.
\bibitem{sep:15} S. R. Ramp et al., ``Preparing to predict: the second autonomous ocean sampling network (AOSN-II) experiment in the Monterey Bay,'' \textit{Deep-Sea Res. II}, vol. 56, no. 3-5, pp. 68--86, 2009.
\bibitem{sep:16} G. Haller and G. Yuan, ``Lagrangian coherent structures and mixing in two-dimensional turbulence,'' \textit{Physica D}, vol. 147, no. 3-4, pp. 352--370, 2000.
\bibitem{sep:17} S. C. Shadden, F. Lekien, and J. E. Marsden, ``Definition and properties of Lagrangian coherent structures from finite-time Lyapunov exponents in two-dimensional aperiodic flows,'' \textit{Physica D}, vol. 212, no. 3-4, pp. 271--304, 2005.
\bibitem{sep:18} G. Haller, ``Lagrangian coherent structures,'' \textit{Annu. Rev. Fluid Mech.}, vol. 47, pp. 137--162, 2015.
\bibitem{sep:19} M. Serra and G. Haller, ``Objective Eulerian coherent structures,'' \textit{Chaos}, vol. 26, no. 5, p. 053110, 2016.
\bibitem{sep:20} P. J. Nolan, M. Serra, and S. D. Ross, ``Finite-time Lyapunov exponents in the instantaneous limit and material transport,'' \textit{Nonlinear Dyn.}, vol. 100, no. 4, pp. 3825--3852, 2020.
\bibitem{sep:21} T. D. Harms, S. L. Brunton, and B. J. McKeon, ``Lagrangian gradient regression for the detection of coherent structures from sparse trajectory data,'' \textit{R. Soc. Open Sci.}, vol. 11, no. 9, p. 240586, 2024.
\bibitem{sep:22} J. Weiss, ``The dynamics of enstrophy transfer in two-dimensional hydrodynamics,'' \textit{Physica D}, vol. 48, no. 2-3, pp. 273--294, 1991.
\bibitem{sep:23} J. C. R. Hunt, A. A. Wray, and P. Moin, ``Eddies, streams, and convergence zones in turbulent flows,'' in \textit{Proc. Summer Program, Center Turbul. Res.}, Rep. CTR-S88, pp. 193--208, 1988.
\bibitem{sep:24} B.-L. Hua and P. Klein, ``An exact criterion for the stirring properties of nearly two-dimensional turbulence,'' \textit{Physica D}, vol. 113, no. 1, pp. 98--110, 1998.
\bibitem{sep:25} C. Dong et al., ``Oceanic mesoscale eddies,'' \textit{Ocean-Land-Atmos. Res.}, vol. 4, p. 0081, 2025.
\bibitem{sep:26} M. Michini, M. A. Hsieh, E. Forgoston, and I. B. Schwartz, ``Robotic tracking of coherent structures in flows,'' \textit{IEEE Trans. Robot.}, vol. 30, no. 3, pp. 593--603, 2014.
\bibitem{sep:27} M. Michini, H. Rastgoftar, M. A. Hsieh, and S. Jayasuriya, ``Distributed formation control for collaborative tracking of manifolds in flows,'' in \textit{Proc. Amer. Control Conf.}, pp. 3874--3880, 2014.
\bibitem{sep:28} M. Michini, M. A. Hsieh, E. Forgoston, and I. B. Schwartz, ``Experimental validation of robotic manifold tracking in gyre-like flows,'' in \textit{Proc. IEEE/RSJ Int. Conf. Intell. Robots Syst.}, pp. 2306--2311, 2014.
\bibitem{sep:29} D. Kularatne and M. A. Hsieh, ``Tracking attracting manifolds in flows,'' \textit{Auton. Robots}, vol. 41, no. 8, pp. 1575--1588, 2017.
\bibitem{sep:30} C. A. Kitts and I. Mas, ``Cluster space specification and control of mobile multirobot systems,'' \textit{IEEE/ASME Trans. Mechatronics}, vol. 14, no. 2, pp. 207--218, 2009.
\bibitem{sep:31} S. Tomer et al., ``A low-cost indoor testbed for multirobot adaptive navigation research,'' in \textit{Proc. IEEE Aerosp. Conf.}, pp. 1--12, 2018.
\bibitem{sep:32} C. Waight and C. A. Kitts, ``A functional indoor testbed for multirobot adaptive navigation in vector field environments,'' in \textit{Proc. ASME Int. Des. Eng. Tech. Conf.}, vol. 89251, 2025.
\bibitem{sep:33} D. Eberly et al., ``Ridges for image analysis,'' \textit{J. Math. Imaging Vis.}, vol. 4, no. 4, pp. 353--373, 1994.
\bibitem{sep:34} NOAA Nat. Centers Environ. Inf., ``Near-real-time surface ocean velocities derived from HF-radar stations,'' NCEI Accession gov.noaa.nodc:IOOS-HFRadarRTVector.
\bibitem{sep:35} I. Mas and C. A. Kitts, ``Obstacle avoidance policies for cluster space control of nonholonomic multirobot systems,'' \textit{IEEE/ASME Trans. Mechatronics}, vol. 17, no. 6, pp. 1068--1079, 2012.
